\documentclass[11pt]{article}
\usepackage[margin=1in]{geometry}
\usepackage{amsmath,amssymb,mathtools}
\usepackage{microtype}
\usepackage[hidelinks]{hyperref}

\newcommand{\R}{\mathbb R}
\newcommand{\Wdot}{\dot W}
\newcommand{\norm}[1]{\lVert #1\rVert}

\title{Lindberg's $L^p$ Jacobian Norm Equivalence\
Was Answered Affirmatively by Hyt\"onen}
\author{Literature-resolution packet for arXiv:1611.02223}
\date{11 August 2026}

\begin{document}
\maketitle

\begin{center}
\fbox{\parbox{0.91\textwidth}{
\textbf{Classification: literature already answered.}
Equation~(7.2) of Lindberg (2017) is true for every dimension $n\ge2$ and
every $1<p<\infty$.  It follows directly from Hyt\"onen's Theorem~3.1.4
(2021 version of arXiv:1804.11167); the specialization is also written out
in Hyt\"onen's proof of Theorem~3.2.1.  This answers the norm-equivalence
subproblem, not Iwaniec's stronger conjecture that every $L^p$ datum is one
Jacobian or admits a continuous choice of preimage.
}}
\end{center}

\section{The exact source question}

In Section~7.3 of \cite{Lindberg}, after proving a first-category result for
inhomogeneous Sobolev maps, Lindberg asks whether
\begin{equation}\label{eq:lindberg}
 \norm{b}_{L^{p'}(\R^n)}
 \asymp
 \sup_{\norm{Du}_{L^{np}(\R^n)}\le1}
 \left|\int_{\R^n} b\,Ju\right|
 \qquad (b\in L^{p'}(\R^n))
\end{equation}
holds for $n\ge2$ and $1<p<\infty$.  This is equation~(7.2) on PDF page 27
(published page 739).  The source calls it an important step toward
Conjecture~7.1, relates its planar case to an Ahlfors--Beurling commutator,
and ends: ``The author has been unable to prove or disprove (7.2).''

The normalization in \eqref{eq:lindberg} uses any fixed finite-dimensional
matrix norm for $Du$; changing that norm changes only the implicit constants.

\section{The later theorem}

Let $d\ge2$.  Hyt\"onen's Theorem~3.1.4 \cite{Hytonen} considers exponents
$r_i\in(1,\infty)$, defined by
\[
 \frac1r=\sum_{i=1}^d\frac1{r_i},
\]
and the quantity
\begin{equation}\label{eq:hytonen}
 \Gamma(b)=
 \sup\left\{
  \left|\int_{\R^d}bJ(u)\right|:
  u=(u_i)_{i=1}^d\in C_c^\infty(\R^d)^d,
  \norm{\nabla u_i}_{L^{r_i}}\le1\ \text{for all }i
 \right\}.
\end{equation}
For $r>1$, the theorem says that $\Gamma(b)<\infty$ exactly when
$b=a+c$ with $a\in L^{r'}$ and $c$ constant, and that the corresponding
function-space norm is comparable to $\Gamma(b)$.

Now take
\[
 r_1=\cdots=r_d=pd.
\]
Then $r=p$.  If $b\in L^{p'}(\R^d)$, the additive constant must be zero, so
Theorem~3.1.4 gives
\begin{equation}\label{eq:specialized}
 \norm{b}_{L^{p'}(\R^d)}\asymp \Gamma(b).
\end{equation}
The constraint in \eqref{eq:hytonen} is equivalent, up to a constant
depending only on $d,p$, to $\norm{Du}_{L^{pd}}\le1$.  Density of
$C_c^\infty$ in the homogeneous Sobolev seminorm and continuity of the
Jacobian pairing identify \eqref{eq:specialized} with Lindberg's
\eqref{eq:lindberg}.

This is not merely an inferred specialization.  On PDF pages 29--30,
Hyt\"onen's proof of Theorem~3.2.1 explicitly writes the dual lower bound
against normalized Jacobians and says that it is ``precisely''
Theorem~3.1.4 with $r=p$ and $r_1=\cdots=r_d=pd$.  Remark~3.2.3(2) then
cites Lindberg's page 739, explains that Lindberg sketched the planar case,
and states that the paper proved the more general result directly with the
Jacobian.

\section{Verdict and scope boundary}

Thus equation~(7.2) has an affirmative answer in all dimensions and for the
full exponent range asked in the source.  Hyt\"onen also proves in
Theorem~3.2.1 that every $L^p$ function is an absolutely controlled
convergent series of Jacobians, i.e. the appropriate weak factorization.

The result does \emph{not} show that every $f\in L^p$ is a single Jacobian.
It therefore does not settle bare surjectivity of
\[
 J:\Wdot^{1,pd}(\R^d,\R^d)\longrightarrow L^p(\R^d),
\]
and it does not construct a continuous right inverse.  Those are the
stronger assertions in Conjecture~7.1.

\paragraph{Evidence locations.}
Source: arXiv:1611.02223, PDF page 27, equation~(7.2) and the final sentence
before the acknowledgements.  Supporting paper: arXiv:1804.11167v4,
Theorem~3.1.4 on PDF page 25; proof of Theorem~3.2.1 on pages 29--30; and
Remark~3.2.3(2) on page 30.

\paragraph{Search and confidence.}
The match is explicit: the later paper cites Lindberg's published page 739
and discusses the same Jacobian/commutator route.  Classification confidence
is high.  The only care needed is the scope boundary between a norming/weak-
factorization result and representation by one Jacobian.

\begin{thebibliography}{9}
\bibitem{Lindberg}
S.~Lindberg,
\emph{On the Hardy Space Theory of Compensated Compactness Quantities},
Arch. Ration. Mech. Anal. \textbf{224} (2017), 709--742;
arXiv:1611.02223.

\bibitem{Hytonen}
T.~P. Hyt\"onen,
\emph{The $L^p$-to-$L^q$ Boundedness of Commutators with Applications to the
Jacobian Operator},
J. Math. Pures Appl. \textbf{156} (2021), 351--391;
arXiv:1804.11167v4.
\end{thebibliography}

\end{document}
